\section{Introduction}

Non-convex optimization algorithms are one of the fundamental tools in modern machine learning, as training neural network models requires optimizing a non-convex loss function. This paper provides a new theoretical framework for building such algorithms. The simplest application of this framework almost exactly recapitulates the standard algorithm used in practice: stochastic gradient descent with momentum (SGDM).
% . Motivated by this, there have been abundant works focused on understanding non-convex optimization from both theoretical and empirical perspectives. 

The goal of any optimization algorithm used to train a neural network is to minimize a potentially non-convex objective function. Formally, given $F:\RR^d\to\RR$, the problem is to solve
\begin{equation*}
    \min_{\vecx\in\RR^d} F(\vecx) = \Ex_{z}[f(\vecx,z)],
\end{equation*}
where $f$ is a stochastic estimator of $F$. In practice, $\vecx$ denotes the parameters of a neural network model, and $z$ denotes the data point. Following the majority of the literature, we focus on first-order stochastic optimization algorithms, which can only access to the stochastic gradient $\nabla f(\vecx,z)$ as an estimate of the unknown true gradient $\nabla F(\vecx)$. We measure the ``cost'' of an algorithm by counting the number of stochastic gradient evaluations it requires to achieve some desired convergence guarantee. We will frequently refer to this count as the number of ``iterations'' employed by the algorithm.

When the objective function is non-convex, finding a global minimum can be intractable. To navigate this complexity, many prior works have imposed various smoothness assumptions on the objective. These include, but not limited to, first-order smoothness \cite{ghadimi2013stochastic, carmon2017convex, arjevani2022lower, carmon2019lower}, second-order smoothness \cite{tripuraneni2018stochastic, carmon2018accelerated, fang2019sharp, arjevani2020second}, and mean-square smoothness \cite{allenzhu2018natasha, fang2018spider, cutkosky2019momentum, arjevani2022lower}. Instead of finding the global minimum, the smoothness conditions allow us to find an $\epsilon$-stationary point $\vecx$ of $F$ such that $\|\nabla F(\vecx)\|\le\epsilon$.

The optimal rates for smooth non-convex optimization are now well-understood. When the objective is smooth, stochastic gradient descent (SGD) requires $O(\epsilon^{-4})$ iterations to find $\epsilon$-stationary point, matching the optimal rate \cite{arjevani2019lower}. When $F$ is second-order smooth, a variant of SGD augmented with occasional random perturbations achieves the optimal rate $O(\epsilon^{-7/2})$ \cite{fang2019sharp, arjevani2020second}. Moreover, when $F$ is mean-square smooth, variance-reduction algorithms, such as SPIDER \cite{fang2018spider} and SNVRG \cite{zhou2018stochastic}, achieve the optimal rate $O(\epsilon^{-3})$ \cite{arjevani2019lower}. All of these algorithms can be viewed as variants of SGD.

In addition to these theoretical optimality results, SGD and its variants are also incredibly effective in practice across a wide variety of deep learning tasks. Among these variants, the family of momentum algorithms have become particularly popular \cite{sutskever2013importance, kingma2014adam, you2017scaling, you2019large, cutkosky2019momentum, cutkosky2020momentum, ziyin2020laprop}. Under smoothness conditions, the momentum algorithms also achieve the same optimal rates.
% ELABORATE: how to draw connections between momentum algorithms and the goal of our paper?

However, modern deep learning models frequently incorporate a range of non-smooth architectures, including elements like ReLU, max pooling, and quantization layers. These components result in a non-smooth optimization objective, violating the fundamental assumption of a vast majority of prior works. Non-smooth optimization is fundamentally more difficult than its smooth counterpart, as in  the worst-case \citet{kornowski2022oracle} show that it is actually  impossible to find a neighborhood around $\epsilon$-stationary. This underscores the need for a tractable convergence criterion in non-smooth non-convex optimization.

One line of research in non-smooth non-convex optimization studies weakly-convex objectives \citep{davis2019stochastic, mai2020convergence}, with a focus on finding $\epsilon$-stationary points of the Moreau envelope of the objectives. It has been demonstrated that various algorithms, including the proximal subgradient method and SGDM, can achieve the optimal rate of $O(\epsilon^{-4})$ for finding an $\epsilon$-stationary point of the Moreau envelope. However, it is important to note that the assumption of weak convexity is crucial for the convergence notion involving the Moreau envelope. Our interest lies in solving non-smooth non-convex optimization without relying on the weak convexity assumption.

To this end, \citet{zhang2020complexity} proposed employing Goldstein stationary points \citep{goldstein1977optimization} as a convergence target in non-smooth non-convex (and non-weakly-convex) optimization. This approach has been widely accepted by recent works studying non-smooth optimization \cite{kornowski2022complexity, lin2022gradient, kornowski2023algorithm, cutkosky2023optimal}. Formally, $\vecx$ is a $(\delta,\epsilon)$-Goldstein stationary point if there exists a random vector $\vecy$ such that $\E[\vecy]=\vecx$, $\|\vecy-\vecx\| \le \delta$ almost surely, and $\|\E[\nabla F(\vecy)]\|\le\epsilon$.\footnote{To be consistent with our proposed definition, we choose to present the definition of $(\delta,\epsilon)$-Goldstein stationary point involving a random vector $\vecy$. This presentation is equivalent to the original definition proposed by \cite{zhang2020complexity}.} The best-possible rate for finding a $(\delta,\epsilon)$-Goldstein stationary point is $O(\delta^{-1}\epsilon^{-3})$ iterations. This rate was only recently achieved  by \citet{cutkosky2023optimal}, who developed an ``online-to-non-convex conversion'' (O2NC) technique that converts online convex optimization (OCO) algorithms to non-smooth non-convex stochastic optimization algorithms. Building on this background, we will relax the definition of stationarity and extend this O2NC technique to eventually develop a convergence analysis of SGDM in the non-smooth and non-convex setting.
% Notably, applied with online gradient descent, O2NC requires $O(\delta^{-1}\epsilon^{-3})$ for finding $(\delta,\epsilon)$-stationary point, which is the optimal rate.


\subsection{Our Contribution}

In this paper, we introduce a new notion of stationarity for non-smooth non-convex objectives. Our notion is a natural relaxation of the Goldstein stationary point, but will allow for  more flexible algorithm design. Intuitively, the problem with the Goldstein stationary point is that to verify that a point $\vecx$ is a stationary point, one must evaluate the gradient many times inside a ball of some small radius $\delta$ about $\vecx$. This means that algorithms finding such points usually make fairly conservative updates to sufficiently explore this ball: in essence, they work by verifying each iterate is \emph{not} close to a stationary point before moving on to the next iterate. Algorithms used in practice do not typically behave this way, and our relaxed definition will not require us to employ such behavior.

% $(c,\epsilon)$-stationary point as an alternative convergence criterion for non-smooth optimization. Formally, we say $\vecx$ is a $(c,\epsilon)$-stationary point if there exists a random vector $\vecy$ such that $\E[\vecy]=\vecx$, $\E\|\vecy-\vecx\|^2\le \epsilon/c$, and $\|\E\nabla F(\vecy)\|\le \epsilon$. With $c=\epsilon/\delta^2$, this definition aligns closely with the standard $(\delta,\epsilon)$-stationary point except for a key modification: it relaxes the deterministic bound of $\vecy$, namely $\|\vecy-\vecx\|\le \delta$, to a probabilistic bound of its variance, $\E\|\vecy-\vecx\|^2 \le \delta^2$. This adjustment expands the scope of algorithms under consideration, accommodating those that may violate the deterministic bound but satisfy the variance bound. Furthermore, similar to $(\delta,\epsilon)$-stationary point, our $(c,\epsilon)$-stationary point criterion automatically reduces to $\epsilon$-stationary point with proper choices of $c$ when the objective is smooth or second-order smooth.

Using our new criterion, we propose a general framework, ``\alg'', that converts OCO algorithms to non-smooth optimization algorithms. This framework is an extension of the  O2NC technique of \citet{cutkosky2023optimal} that distinguishes itself through two significant improvements.

Firstly, the original O2NC method requires the OCO algorithm to constrain all of its iterates to a small ball of radius roughly $\delta\epsilon^2$. 
% to periodically restarting the online learning algorithm every $T$ iterations (where typically $T=\sqrt{N}$ with $N$ denoting the total number of iterations). 
This approach is designed to ensure that the parameters within any period of $\epsilon^{-2}$ iterations remain inside a ball of radius $\delta$. The algorithm then uses these $\epsilon^{-2}$ gradient evaluations inside a ball of radius $\delta$ to check if the current iterate is a stationary point (i.e., if the average gradient has norm less than $\epsilon$). Our new criterion, however, obviates the need for such explicit constraints, intuitively allowing our algorithms to make larger updates when far from a stationary point.
% We demonstrate that, even without periodic resets, it is possible for any given iterate $\vecx_n$ to be associated with a random vector $\vecy_n$ whose variance remains controlled. This allows the OCO algorithm to proceed uninterrupted.

Secondly, O2NC does not evaluate gradients at the actual iterates. Instead, gradients are evaluated at an intermediate variable $\vecw_n$ lying between the two iterates $\vecx_n$ and $\vecx_{n+1}$. This conflicts with essentially all practical  algorithms, and moreover imposes an extra memory burden. In contrast, our algorithm evaluates  gradients exactly at each iterate, which simplifies implementation and improves space complexity.
% by setting $s_n$ to an exponential random variable, our algorithm eliminates the requirement for this intermediate state, further improving the memory complexity.

Armed  with this improved framework, we proposed an unconstrained variant of online gradient descent, which is derived from the family of online mirror descent with composite loss. When applied within this algorithm, our framework produces an algorithm that is exactly equal to stochastic gradient descent with momentum (SGDM), subject to an additional random scaling on the update. Notably, it also achieves the optimal rate under our new criterion.
% Next, we consider a standard adaptive learning rate schedule for online gradient descent. Integrated with \alg, this approach evolves into an adaptive momentum algorithm. 

To summarize, this paper has the following contributions:
\begin{itemize}
    \item We introduce a relaxed convergence criterion for non-smooth optimization that recovers all useful properties of Goldstein stationary point.

    \item We propose a modified online-to-non-convex conversion framework that does not require intermediate states.

    \item We apply our new conversion to the most standard OCO algorithm: ``online gradient descent''. The resulting method achieves optimal convergence guarantees as is almost exactly the same as the standard SGDM algorithm. The only difference is that the updates of  SGDM are now scaled by an exponential random variable. This is especially remarkable because the machinery that we employ does not particularly resemble SGDM until it is finally all put together.
    % Further applied with adaptive learning rate, our algorithm achieves adaptive convergence guarantee.
    
    % our algorithm finds a $(c,\epsilon)$-stationary point within $O(c^{1/2}\epsilon^{-7/2})$ iterations. With $c=\epsilon/\delta^2$, this recovers the optimal rate $O(\delta^{-1}\epsilon^{-3})$ under the standard $(\delta,\epsilon)$-stationary point criterion. Notably, this algorithm recovers SGDM, thus providing a theoretical proof for its empirical success.

    % \item Applied with adaptive learning rate, our algorithm again achieves the $O(c^{1/2}\epsilon^{-7/2})$ rate while recovering an adaptive momentum algorithm.
\end{itemize}